\documentclass[a4paper, 14pt]{extarticle}

\usepackage{../latexDependencies/misc/preamble}

\geometry{a4paper}

% Название дисциплины
\newcommand{\subject}{Теория вероятности и математическая статистика} 

% Тип работы
% lab - для лабораторной работы 
% hw  - для домашней     работы
\newcommand{\task}{lab} 

% Номер работы
\newcommand{\taskNumber}{1} 

% Название работы
\newcommand{\taskName}{Первоначальная обработка статистических данных} 

% Имя студента
\newcommand{\studentName}{Очкин Н.В.}

% Имя преподававателя
\newcommand{\teacherName}{Облакова Т.В.}

% Группа
\newcommand{\group}{ФН11-52Б}

% Вариант
\newcommand{\variant}{9}

\begin{document}

\graphicspath{ {../latexDependencies/images} } 
\normalsize

\newcommand{\printTask}{%
    \ifthenelse{\equal{\task}{lab}}{%
        лабораторной%
    }{%
        \ifthenelse{\equal{\task}{hw}}{%
            домашней%
        }{%
            Неизвестный тип задания%
        }%
    }%
}

\begin{titlepage}

    \begin{center}
        {\footnotesize \itshape Федеральное государственное бюджетное 
                       образовательное учреждение высшего образования}
    \end{center}

    \begin{minipage}[c]{0.1\textwidth}
        \includegraphics[width=1.1\textwidth]{iconBMSTU}
    \end{minipage}
    \hfill
    \begin{minipage}[c]{0.9\textwidth}
        \centering
        \itshape
        \bfseries
        \small
        \guillemotleft Московский государственный технический университет \\
        имени Н.Э. Баумана\guillemotright \\
        (национальный исследовательский университет) \\
        (МГТУ им. Н.Э. Баумана) 
    \end{minipage}

    \vspace{0.5cm}
    \noindent\rule{\textwidth}{2pt} \\

    \noindent\uline{\textbf{ФАКУЛЬТЕТ} ФУНДАМЕНТАЛЬНЫЕ НАУКИ} \\
    \vspace{-5pt} \\
    \noindent\uline{\textbf{КАФЕДРА} ВЫЧИСЛИТЕЛЬНАЯ МАТЕМАТИКА И МАТЕМАТИЧЕСКАЯ} \\
    \vspace{-5pt} \\
    \noindent\uline{ФИЗИКА (ФН11)} \\
    \vspace{-5pt} \\
    \noindent\uline{\textbf{НАПРАВЛЕНИЕ ПОДГОТОВКИ} МАТЕМАТИКА И КОМПЬЮТЕРНЫЕ} \\
    \vspace{-5pt} \\
    \noindent\uline{НАУКИ (02.03.01)} \\

    \begin{center}
        \bfseries
        \textsc{О т ч е т} \\[10pt]
        по \printTask {} работе \textnumero {} \taskNumber
    \end{center}

    \vspace{10pt}

    \hspace{10pt} 
    \noindent \textbf{Название \printTask {} работы:} \par
    \vspace{5pt}
    \hspace{10pt} 
    \noindent \textbf{\uline{\taskName}}

    \vspace{10pt}

    \begin{center}
        \bfseries
        Вариант \textnumero {} \variant
    \end{center}

    \vspace{20pt}

    \hspace{10pt} 
    \noindent \textbf{Дисциплина:} \par
    \vspace{5pt}
    \hspace{10pt} 
    \noindent \subject

    \vspace{10pt}

    \begin{flushright}
        \renewcommand{\arraystretch}{3}
        \begin{tabular}{r r r}
            \multicolumn{1}{l}{Студент группы \uline{\group}} & 
            $\quad \underset{\text{(Подпись, дата)}}{\underline{\hspace{3cm}}} \quad$ & 
            \multicolumn{1}{c}{$\underset{\text{(И.О. Фамилия)}}{\uline{\textbf{\studentName}}}$} \\

            \multicolumn{1}{l}{Преподаватель} & 
            $\quad \underset{\text{(Подпись, дата)}}{\underline{\hspace{3cm}}} \quad$ & 
            \multicolumn{1}{c}{$\underset{\text{(И.О. Фамилия)}}{\uline{\textbf{\teacherName}}}$} \\
        \end{tabular}
    \end{flushright}

    \vfill

    \begin{center}
        \small
        Москва, 2024
    \end{center}
\end{titlepage}


% Reset the font size and page margins
\newgeometry{left=25mm, right=25mm, top=20mm, bottom=20mm}
% \fontsize{12pt}{16pt}\selectfont

\graphicspath{ {../latexDependencies/images/LW1} }

% Customize subsection style
\titleformat{\subsection}
  {\normalfont\normalsize\bfseries}{\thesubsection}{1em}{}

\pagenumbering{gobble}

\tableofcontents
\newpage

\pagenumbering{arabic}
\setcounter{page}{1}

\setstretch{1}
\linespread{1.1}

% --------------------------------------START--------------------------------------

\fontsize{12pt}{16pt}\selectfont

\section{Задание}

По данной выборке

\begin{enumerate}
    \item Найдите крайние члены вариационного ряда и размах выборки
    \item Осуществите группировку данных (количество интервалов 
    находим по правилу Стерджеса) 
    \item По сгруппированным данным постройте гистограмму 
    относительных частот
    \item Вычислите выборочное среднее и выборочную дисперсию.
    \item По виду гистограммы определите возможный закон распределения, 
    оцените параметры этого закона по методу моментов, постройте совмещенные 
    графики гистограммы и плотности предполагаемого закона
    \item Найдите эмпирическую функцию распределения и постройте совмещенные 
    графики эмпирической и теоретической функций распределения
\end{enumerate}

\section{Исходные данные}

\begin{center}
    \begin{table}[h!]
        \begin{adjustbox}{max width=\textwidth}
            \begin{tabular}{|c|c|c|c|c|c|c|c|c|c|c|}
                \hline
                14.495 & 4.715 & 7.175 & 8.428 & 11.093 &
                3.375 & 12.906 & 8.415 & 8.916 & 13.48 \\
                \hline
                5.343 & 17.985 & 15.992 & 13.89 & 9.838 &
                13.924 & 9.012 & 9.458 & 17.69 & 6.542 \\
                \hline
                14.396 & 8.592 & 8.206 & 14.237 & 7.357 &
                10.821 & 12.767 & 16.058 & 12.959 & 4.354 \\
                \hline
                12.888 & 10.268 & 9.182 & 5.647 & 8.282 &
                2.903 & 15.988 & 12.959 & 14.919 & 6.339 \\
                \hline
                2.375 & 17.921 & 9.097 & 15.85 & 11.449 &
                11.095 & 9.493 & 12.175 & 7.479 & 13.535 \\
                \hline
                9.234 & 6.078 & 4.964 & 6.355 & 13.957 &
                12.911 & 15.694 & 14.286 & 9.869 & 5.175 \\
                \hline
                5.811 & 7.241 & 5.814 & 3.086 & 6.875 &
                3.878 & 5.333 & 15.134 & 12.924 & 9.159 \\
                \hline
                4.727 & 4.646 & 15.535 & 9.919 & 17.117 &
                10.351 & 16.892 & 12.423 & 10.511 & 4.942 \\
                \hline
                4.843 & 9.927 & 15.864 & 3.635 & 17.963 &
                8.25 & 5.14 & 6.734 & 12.622 & 13.325 \\
                \hline
                3.377 & 16.195 & 12.04 & 12.768 & 2.744 &
                14.186 & 9.354 & 15.439 & 14.612 & 15.649 \\
                \hline
                8.681 & 5.006 & 3.608 & 2.867 & 12.177 &
                15.506 & 7.683 & 14.022 & 17.103 & 8.905 \\
                \hline
                12.173 & 17.757 & 6.883 & 2.666 & 9.861 &
                5.743 & 16.175 & 15.308 & 7.039 & 15.238 \\
                \hline
            \end{tabular}
        \end{adjustbox}
        \caption{Исходные данные}
        \label{tab:data}
    \end{table}
\end{center}

\newpage

\section{Решение}
Данная работа была выполнена с использованием языка программирования 
\textsc{python} и следующих модулей и библиотек: 

\renewcommand{\labelitemi}{\scalebox{0.5}{$\bullet$}}

\begin{center}
    \begin{multicols}{4}
        \begin{itemize}[itemsep=0pt]
            \item numpy
            \item pandas
            \item matplotlib
            \item scipy
            \item IPython
            \item math
            \item pyplot
            \item display
            \item stats
        \end{itemize}
    \end{multicols}
\end{center}

\subsection{Найдите крайние члены вариационного ряда и размах выборки}

% --------------------------------------CODE---------------------------------------

\section{Приложение}

Программный код, с помощью которого была выполнена данная лабораторная работа.\\

\vspace{20pt}

\noindent \textbf{Примечание}.\\

\noindent Так как отчет был написан с использованием дистрибутива TeX и 
следующий код был отформатирован с использованием окружения \textsc{lstlisting}, 
в некоторых местах текст, написанный на русском языке, может иметь проблемы 
с выравниванием, пробелами, шрифтом и т.д. Это связано с тем, что библиотека 
\textsc{listings}, из которой мы берем окружение для форматирования кода, 
плохо работает с Unicode.

\vspace{20pt}

% \newgeometry{left=15mm, right=15mm, top=20mm, bottom=20mm}

\setstretch{1}

\lstdefinestyle{mystyle}{
    basicstyle=\ttfamily\footnotesize,
    keywordstyle=\color{blue},
    stringstyle=\color{red},
    commentstyle=\color{green!50!black},
    showstringspaces=false
}

% Set the style for Python code
\lstset{style=mystyle, language=Python, extendedchars=\true}

\begin{lstlisting}
import numpy as np
import pandas as pd
import matplotlib.pyplot as plt
import scipy as sp
import math
from IPython.display import display, Math

PRECISION = 5
pd.set_option('display.float_format', '{:.3f}'.format)
np.set_printoptions(precision=PRECISION)
\end{lstlisting}

\begin{lstlisting}
_data = [
    14.495, 4.715,  7.175,  8.428,  11.093, 3.375
    12.906, 8.415,  8.916,  13.48,  5.343,  17.985
    15.992, 13.89,  9.838,  13.924, 9.012,  9.458
    17.69,  6.542,  14.396, 8.592,  8.206,  14.237
    7.357,  10.821, 12.767, 16.058, 12.959, 4.354
    12.888, 10.268, 9.182,  5.647,  8.282,  2.903
    15.988, 12.959, 14.919, 6.339,  2.375,  17.921
    9.097,  15.85,  11.449, 11.095, 9.493,  12.175
    7.479,  13.535, 9.234,  6.078,  4.964,  6.355
    13.957, 12.911, 15.694, 14.286, 9.869,  5.175
    5.811,  7.241,  5.814,  3.086,  6.875,  3.878
    5.333,  15.134, 12.924, 9.159,  4.727,  4.646
    15.535, 9.919,  17.117, 10.351, 16.892, 12.423
    10.511, 4.942,  4.843,  9.927,  15.864, 3.635
    17.963, 8.25,   5.14,   6.734,  12.622, 13.325
    3.377,  16.195, 12.04,  12.768, 2.744,  14.186
    9.354,  15.439, 14.612, 15.649, 8.681,  5.006
    3.608,  2.867,  12.177, 15.506, 7.683,  14.022
    17.103, 8.905,  12.173, 17.757, 6.883,  2.666
    9.861,  5.743,  16.175, 15.308, 7.039,  15.238
]
\end{lstlisting}

\begin{lstlisting}
# Variation series
_parsedData = [_data[i:i+10] for i in range(0, len(_data), 10)] 

_parsedData
\end{lstlisting}

\begin{center}
    1. Крайние члены вариационного ряда и размах выборки
\end{center}

\begin{lstlisting}
# Extreme members
_maxEl = np.max(_parsedData)
_minEl = np.min(_parsedData)

print(f"Крайние члены: {_maxEl}, {_minEl}")

# Sample range
_range = np.round(_maxEl - _minEl, PRECISION)

print(f"Размах выборки: {_range}")
\end{lstlisting}

\begin{center}
    2. Осуществите группировку данных
\end{center}

\begin{lstlisting}
# Let's find the number of intervals using Sturges' rule
_n = len(_data)
_k = 1 + np.trunc(np.log2(_n))

print(f"Количество интервалов: {_k}")

# Interval step
_h = _range / _k

print(f"Шаг интервала: {_h}")
\end{lstlisting}

\begin{lstlisting}
# function for finding frequencies
def findFrequency(minim, maxim, data, islast = False):
    count = 0

    for el in data:
        if minim <= el < maxim:
            count += 1

    return count
\end{lstlisting}
    
\begin{lstlisting}
intervals           = [i for i in range(1, int(_k) + 1)] # interval number
tableRanges         = [] # interval
frequencies         = [] # frequency
relativeFrequencies = [] # relative frequency
midRanges           = [] # middle of the interval

curr = _minEl
while (math.floor(curr + _h) <= _maxEl):
    firstEl  = curr                        # first element of the current interval 
    secondEl = round(curr + _h, PRECISION) # second element of the current interval
    tableRanges.append((firstEl, secondEl))

    midRange = round((firstEl + secondEl)/2, PRECISION) # middle of the interval
    midRanges.append(midRange)

    frequency = findFrequency(firstEl, secondEl, _data) \
        if not (abs(curr + _h - _maxEl) < 0.1)          \
        else findFrequency(firstEl, secondEl, _data, True) # frequency
    frequencies.append(frequency)                          
    
    relativeFrequency = round(frequency / _n, PRECISION) # relative frequency
    relativeFrequencies.append(relativeFrequency)

    curr = round(curr + _h, PRECISION)
\end{lstlisting}

\begin{lstlisting}
_intervalRange = pd.DataFrame({'номер интервала':       intervals, 
                               'интервал':             tableRanges, 
                               'середина интервала':     midRanges,
                               'частота':              frequencies,
                               'относительная частота': relativeFrequencies})

_intervalRange
\end{lstlisting}

\begin{lstlisting}
# interval boundaries
_int1 = midRanges - _h/2
_int1 = np.append(_int1, _maxEl)

print(f'границы интервалов: {_int1}')
\end{lstlisting}

\begin{center}
    3. По сгруппированным данным постройте гистограмму относительных частот
\end{center}

\begin{lstlisting}
def buildLine(x, y, color):
    plt.plot(x, y, color=color, linestyle='-', linewidth=1.5)
\end{lstlisting}

\begin{lstlisting}
def buildBar(x, y, overlay = None, showAxis = True):
    RED   = '#6F1D1B'
    BLUE  = '#00ADB5'
    WHITE = '#EEEEEE'
    BLACK = '#393E46' # #253031 #0D1B1E #1C2321

    # Define font sizes
    SIZE_DEFAULT = 14
    SIZE_LARGE   = 16
    SIZE_TICKS   = 10
    plt.rc("font", weight="normal")        # controls default font
    plt.rc("font", size=SIZE_DEFAULT)      # controls default text sizes
    plt.rc("axes", titlesize=SIZE_LARGE)   # fontsize of the axes title
    plt.rc("axes", labelsize=SIZE_DEFAULT) # fontsize of the x and y labels
    plt.rc("xtick", labelsize=SIZE_DEFAULT)# fontsize of the tick labels
    plt.rc("ytick", labelsize=SIZE_DEFAULT)# fontsize of the tick labels

    fig, ax = plt.subplots(
        figsize=(6, 5)
    )

    xticks = [i for i in range(math.floor(min(x)) - 1, math.ceil(max(x)) + 2)]

    # Hide the all but the bottom spines (axis lines)
    ax.spines["right"].set_visible(False)
    ax.spines["left"].set_visible(False)
    ax.spines["top"].set_visible(False)

    if not showAxis:
        ax.spines["bottom"].set_visible(False)

    # Only show ticks on the left and bottom spines
    ax.yaxis.set_ticks_position("left")
    ax.xaxis.set_ticks_position("bottom")
    ax.spines["bottom"].set_bounds(min(xticks), max(xticks))

    # histogtamm 
    plt.bar(x, y, width=2.25, color='white', edgecolor='black', 
            linewidth=1.5, align='center')

    # line
    if overlay is None:
        buildLine(x, y, RED)
    else:
        xmin, xmax = plt.xlim()
        x = np.linspace(xmin, xmax, 100)
        calculate_y = overlay(x)
        plt.plot(x, calculate_y, color=BLUE, linestyle='-', linewidth=1.5)

    # axis names
    if showAxis:
        plt.xlabel('int')
        plt.ylabel('$\\frac{p_k}{h}$', fontsize=20)
        plt.xticks(xticks)
    else:
        ax.xaxis.set_ticks_position('none')
        ax.yaxis.set_ticks_position('none')
        ax.set_xticks([])
        ax.set_yticks([])

    # Adjust the font size of the tick labels
    plt.tick_params(axis='both', which='major', labelsize=SIZE_TICKS)

    plt.show()
\end{lstlisting}

\begin{lstlisting}
def overlayPlot(overlay):
    BLUE  = '#00ADB5'

    plt.hist(_data, bins=int(_k), color='white', edgecolor='black', 
             density=True, alpha=0.6)

    xmin, xmax = plt.xlim()
    x = np.linspace(xmin, xmax, 100)
    
    calculate_y, label = overlay(x)
    plt.plot(x, calculate_y, color=BLUE, linestyle='-', linewidth=2, label=label)

    plt.legend()
    plt.show()
\end{lstlisting}

\begin{lstlisting}
# is indicated along the ordinate axis when constructing a histogram
_histogrammOrdinateAxis = relativeFrequencies / _h

print(f'x: {midRanges}')
print(f'y: {_histogrammOrdinateAxis}')

buildBar(midRanges, _histogrammOrdinateAxis)
\end{lstlisting}

\begin{lstlisting}
plt.hist(_data, bins=int(_k))
plt.show()
\end{lstlisting}

\begin{center}
    4. Вычислите выборочное среднее и выборочную дисперсию
\end{center}

\begin{lstlisting}
# sample average
_overlineX = (1 / _n) * np.sum(_data) 

# sample variance
_S2 = 1 / (_n - 1) * np.sum((_data - _overlineX)**2)

print(f'выборочное среднее: {_overlineX}')
print(f'выборочная дисперсия: {_S2}')
\end{lstlisting}

\begin{center}
    5. По виду гистограммы определите возможный закон распределения, 
    оцените параметры этого закона по методу моментов, постройте совмещенные 
    графики гистограммы и плотности предполагаемого закона
\end{center}

\begin{lstlisting}
def align(name, enName, kstestRes, *params):
    def line(n):
        display('-'*n)

    heading: str = f"\\begin{{equation}} " \
                   f"   \\text{{{name}}} " \
                   f"\\end{{equation}}   " \
    
    dataDisplay: str = f"\\begin{{align}}\n" 

    NAME   = 0
    SYMBOL = 1
    VALUE  = 2
    for param in params:
        dataDisplay += f"& \\text{{{param[NAME]}}}, "\
                       f"{param[SYMBOL]}: & \\quad {param[VALUE]} \\\\\n"

    dataDisplay += f"\\end{{align}}" 

    kstestHeading = \
    f"\\begin{{equation}} "                           \
    f"   \\text{{результаты теста Колмогорова-Смирнова}}: " \
    f"\\end{{equation}} "                                  

    statistic          = kstestRes.statistic
    pvalue             = kstestRes.pvalue
    statistic_location = kstestRes.statistic_location
    statistic_sign     = kstestRes.statistic_sign

    statisticDisplay = f'максимальное различие между эмпирической и '      \
                       f'теоретической функциями распределения составляет ' \
                       f'примерно: {int(np.round(statistic, 2) * 100)}\\%'

    pvalueFlag = True if pvalue < 0.05 else False

    pvalueDisplay = f'значение p-уровня значимости: {pvalue} '
    pvalueDisplay += '<' if pvalueFlag else '>'
    pvalueDisplay += f' 0.05'

    if pvalueFlag:
        pvalueDisplay += \
        f' (данные не соответствуют предполагаемому распределению.)'
    else:
        pvalueDisplay += \
        f' (на уровне значимости 0.05 нет оснований отвергнуть гипотезу ' \
        f'о соответствии данных выбранному распределению)'

    statistic_locationDisplay = \
    f'максимальное отклонение наблюдается при значении данных, ' \
    f'равном {statistic_location}'

    statistic_signDisplay = \
    f'в точке максимального отклонения эмпирическая функция распределения '
    statistic_signDisplay += 'ниже' if statistic_sign < 0 else 'выше'
    statistic_signDisplay += ', чем теоретическая'

    kstestResDisplay = [statisticDisplay, 
                        pvalueDisplay, 
                        statistic_locationDisplay, 
                        statistic_signDisplay]

    kstestDisplay  = f"\\begin{{align}}\n" 

    for res in kstestResDisplay:
        kstestDisplay += f"& \\text{{{res}}} \\\\\n"

    kstestDisplay += f"\\end{{align}}" 

    line(200)

    display(Math(heading))
    display(Math(dataDisplay))
    display(Math(kstestHeading))
    display(Math(kstestDisplay))

    match enName:
        case 'norm': 
            def buildLaw(x):
                return sp.stats.norm.pdf(
                            x, 
                            np.mean(_data), 
                            np.std(_data, ddof=1)
                        ), name

        case 'chi2': 
            def buildLaw(x):
                df_chi2 = 2
                return sp.stats.chi2.pdf(
                            x, 
                            df_chi2, 
                            loc=np.min(_data), 
                            scale=np.std(_data)
                        ), name + f'(df={df_chi2})'

        case 'expon': 
            def buildLaw(x):
                return sp.stats.expon.pdf(
                            x, 
                            loc=np.min(_data), 
                            scale=np.mean(_data)-np.min(_data)
                        ), name

        case 'gamma': 
            def buildLaw(x):
                shape_gamma = 2
                return sp.stats.gamma.pdf(
                            x, 
                            shape_gamma, 
                            loc=np.min(_data), 
                            scale=np.std(_data)
                        ), name + f'(shape={shape_gamma})'
            
        case 'poisson': 
            def buildLaw(x):
                lambda_poisson = np.mean(_data) 
                return sp.stats.poisson.pmf(
                            np.floor(x), 
                            lambda_poisson
                        ), name
            
        case 'uniform': 
            def buildLaw(x):
                return sp.stats.uniform.pdf(
                            x, 
                            loc=np.min(_data), 
                            scale=np.ptp(_data)
                        ), name

        case 't': 
            def buildLaw(x):
                df_t = 2
                return sp.stats.t.pdf(
                            x, 
                            df_t, 
                            loc=np.mean(_data), 
                            scale=np.std(_data)
                        ), name + f'(df={np.round(df_t, PRECISION)})'

        case 'lognorm': 
            def buildLaw(x):
                shape_lognorm = np.std(np.log(_data))
                return sp.stats.lognorm.pdf(
                            x, 
                            shape_lognorm, 
                            loc=np.min(_data), 
                            scale=np.exp(np.mean(np.log(_data)))
                        ), name

        case 'beta':  
            def buildLaw(x):
                a_beta = params[0][VALUE]
                b_beta = params[1][VALUE]
                return sp.stats.beta.pdf(
                            (x - np.min(_data)) / np.ptp(_data), 
                            a_beta, 
                            b_beta
                        ) / np.ptp(_data), \
                        name + f'(a={np.round(a_beta, PRECISION)},'\
                               f'\n b={np.round(b_beta, PRECISION)})'
            
        case 'weibull_min': 
            def buildLaw(x):
                shape_weibull = params[0][VALUE]
                return sp.stats.weibull_min.pdf(
                            x, 
                            shape_weibull, 
                            loc=np.min(_data), 
                            scale=np.std(_data)
                        ), \
                        name + f'(shape={np.round(shape_weibull, PRECISION)})'

        case 'pareto': 
            def buildLaw(x):
                b_pareto = 0.1
                return sp.stats.pareto.pdf(
                            x - np.min(_data) + 1, 
                            b_pareto
                        ), name + f'(shape={b_pareto})'
            
        case default:
            pass

    overlayPlot(buildLaw)
\end{lstlisting}

\begin{lstlisting}
def fillArrays(*args):
    allNames              = args[2]
    allStatistic          = args[3]
    allPvalue             = args[4]
    allStatistic_location = args[5]
    allStatistic_sign     = args[6]

    name = args[1]

    statistic          = args[0].statistic
    pvalue             = args[0].pvalue
    statistic_location = args[0].statistic_location
    statistic_sign     = args[0].statistic_sign

    arrays = [allNames, 
              allStatistic, 
              allPvalue, 
              allStatistic_location, 
              allStatistic_sign]
    data   = [name, 
              statistic, 
              pvalue, 
              statistic_location, 
              statistic_sign]

    for ind, el in enumerate(data):
        arrays[ind].append(el)
\end{lstlisting}

\begin{lstlisting}
data = np.array(_data)

allNames              = []
allStatistic          = []
allPvalue             = []
allStatistic_location = []
allStatistic_sign     = []

# Normal distribution (norm)
mean, std = sp.stats.norm.fit(data)
result = sp.stats.kstest(data, 'norm', args=(mean, std))
fillArrays(result, 
            'norm', 
            allNames, 
            allStatistic, 
            allPvalue, 
            allStatistic_location, 
            allStatistic_sign)
align('нормальное', 'norm', result, ['среднее', '\\mu', mean], 
                                    ['стандратное отклонение', '\\sigma', std])

# Chi-square distribution (chi2)
df, loc, scale = sp.stats.chi2.fit(data)
result = sp.stats.kstest(data, 'chi2', args=(df,))
fillArrays(result, 
            'chi2', 
            allNames, 
            allStatistic, 
            allPvalue, 
            allStatistic_location, 
            allStatistic_sign)
align('хи-квадрат', 'chi2', result, ['количество степеней свободы', '', df], 
                                    ['сдвиг ', '', loc],
                                    ['параметр масштаба', '', scale])

# Exponential distribution (expon)
loc, scale = sp.stats.expon.fit(data)
result = sp.stats.kstest(data, 'expon', args=(loc, scale))
fillArrays(result, 
            'expon', 
            allNames, 
            allStatistic, 
            allPvalue, 
            allStatistic_location, 
            allStatistic_sign)
align('экспоненциальное', 'expon', result, 
      ['сдвиг', '', loc], 
      ['параметр масштаба', '\\frac{1}{\\lambda}', scale])

# Gamma distribution (gamma)
shape, loc, scale = sp.stats.gamma.fit(data)
result = sp.stats.kstest(data, 'gamma', args=(shape, loc, scale))
fillArrays(result, 
            'gamma', 
            allNames, 
            allStatistic, 
            allPvalue, 
            allStatistic_location, 
            allStatistic_sign)
align('гамма-распределение', 'gamma', result, ['параметр формы', '', shape], 
                                              ['сдвиг ', '', loc],
                                              ['параметр масштаба', '', scale])

# Poisson distribution (poisson)
lambda_ = np.mean(data)
result = sp.stats.kstest(data, 'poisson', args=(lambda_,))
fillArrays(result, 
            'poisson', 
            allNames, 
            allStatistic, 
            allPvalue, 
            allStatistic_location, 
            allStatistic_sign)
align('пуассоновское', 'poisson', result, ['параметр интенсивности', '\\mu', lambda_])

# Uniform distribution (uniform)
loc, scale = sp.stats.uniform.fit(data)
result = sp.stats.kstest(data, 'uniform', args=(loc, scale))
fillArrays(result, 
            'uniform', 
            allNames, 
            allStatistic, 
            allPvalue, 
            allStatistic_location, 
            allStatistic_sign)
align('равномерное', 'uniform', result, ['нижняя граница', '', loc], 
                                        ['размах', '', scale])

# Student's t-distribution (t)
df, loc, scale = sp.stats.t.fit(data)
result = sp.stats.kstest(data, 't', args=(df, loc, scale))
fillArrays(result, 
            't', 
            allNames, 
            allStatistic, 
            allPvalue, 
            allStatistic_location, 
            allStatistic_sign)
align('t-распределение Стьюдента', 't', result, 
      ['количество степеней свободы', '', df], 
      ['сдвиг', '', loc],
      ['масштаб', '', scale])

# Lognormal distribution (lognorm)
shape, loc, scale = sp.stats.lognorm.fit(data)
result = sp.stats.kstest(data, 'lognorm', args=(shape, loc, scale))
fillArrays(result, 
            'lognorm', 
            allNames, 
            allStatistic, 
            allPvalue, 
            allStatistic_location, 
            allStatistic_sign)
align('Логнормальное', 'lognorm', result, ['параметр формы', '', shape], 
                                          ['сдвиг', '', loc],
                                          ['параметр масштаба', '', scale])

# Beta distribution (beta)
a, b, loc, scale = sp.stats.beta.fit(data)
result = sp.stats.kstest(data, 'beta', args=(a, b, loc, scale))
fillArrays(result, 
            'beta', 
            allNames, 
            allStatistic, 
            allPvalue, 
            allStatistic_location, 
            allStatistic_sign)
align('бета-распределение', 'beta', result, ['параметр формы 1', '', a], 
                                         ['параметр формы 2', '', b],
                                         ['сдвиг', '', loc],
                                         ['масштаб', '', scale])

# Weibull distribution (weibull_min)
c, loc, scale = sp.stats.weibull_min.fit(data)
fillArrays(result, 
            'weibull_min', 
            allNames, 
            allStatistic, 
            allPvalue, 
            allStatistic_location, 
            allStatistic_sign)
align('вейбулловское ', 'weibull_min', result, ['параметр формы', '', c], 
                                               ['сдвиг', '', loc],
                                               ['масштаб', '', scale])

# Pareto distribution (pareto)
b, loc, scale = sp.stats.pareto.fit(data)
result = sp.stats.kstest(data, 'pareto', args=(b, loc, scale))
fillArrays(result, 
            'pareto', 
            allNames, 
            allStatistic, 
            allPvalue, 
            allStatistic_location, 
            allStatistic_sign)
align('парето ', 'pareto', result, ['параметр формы', '', b], 
                                   ['сдвиг', '', loc],
                                   ['масштаб', '', scale])

_kstestData = pd.DataFrame({
    'name': allNames,
    'statistic': allStatistic,
    'pvalue': allPvalue,
    'statistic_location': allStatistic_location,
    'statistic_sign': allStatistic_sign
})

_kstestData
\end{lstlisting}

\begin{lstlisting}
# pvalue - If the pvalue is less than a predetermined significance level 
# (for example, 0.05), then the null hypothesis is rejected, which means 
# that the data does not fit the expected distribution.
_kstestData_filtered = _kstestData[_kstestData['pvalue'] > 0.05]
_kstestData_filtered
\end{lstlisting}

\begin{lstlisting}
# statistic - The larger the value, the greater the difference between the 
# data and the theoretical distribution.
_kstestData_filtered.sort_values(by='statistic')
\end{lstlisting}

\begin{lstlisting}
# Calculate sampling parameters
mean = np.mean(_data)
std = np.std(_data, ddof=1)

# Sample histogram
plt.hist(_data, bins=7, density=True, alpha=0.6, color='g')

xmin, xmax = plt.xlim()
x = np.linspace(xmin, xmax, 100)

# Check for normal distribution
p = sp.stats.norm.pdf(x, mean, std)
plt.plot(x, p, 'k', linewidth=2, label='Normal')

# Check for uniform distribution
uniform_fit = sp.stats.uniform.pdf(
                    x, 
                    loc=np.min(_data), 
                    scale=np.ptp(_data)
              )
plt.plot(x, uniform_fit, 'b', linewidth=2, label='Uniform')

# Check for beta distribution
a_beta, b_beta = 1.2406801559765772, 0.9441200264826078  # Параметры формы
beta_fit = sp.stats.beta.pdf(
                (x - np.min(_data)) / np.ptp(_data), 
                a_beta, 
                b_beta
           ) / np.ptp(_data)
plt.plot(x, beta_fit, 'cyan', linewidth=2, 
                              label='Beta (a=1.24068, b=0.94412)')

# Add a legend and display the graph
plt.legend()
plt.show()
\end{lstlisting}

\begin{lstlisting}
_param = 2 * _overlineX

_param
\end{lstlisting}

\begin{center}
    6. Найдите эмпирическую функцию распределения и постройте совмещенные 
    графики эмпирической и теоретической функций распределения
\end{center}

\begin{lstlisting}
# accumulated frequency column
_kum = np.zeros(np.size(relativeFrequencies) + 1)
ind = 1 
for relativeFrequency in relativeFrequencies:
    _kum[ind] = _kum[ind-1] + relativeFrequencies[ind-1]
    ind += 1
print(f'kum: {_kum}')
\end{lstlisting}

\begin{lstlisting}
def femp(x):
    def ind(x):
        return 1 if x > 0 else 0

    sumy = 0 
    for i in range(int(_k)):
        sumy += relativeFrequencies[i] * ind(x - midRanges[i])
    
    return sumy
\end{lstlisting}

\begin{lstlisting}
def buildFemp(cdf_y_values):
    RED   = '#6F1D1B'
    BLUE  = '#00ADB5'
    WHITE = '#EEEEEE'
    BLACK = '#393E46' # #253031 #0D1B1E #1C2321

    # Define font sizes
    SIZE_DEFAULT = 14
    SIZE_LARGE   = 16
    SIZE_TICKS   = 10
    plt.rc("font", weight="normal")        # controls default font
    plt.rc("font", size=SIZE_DEFAULT)      # controls default text sizes
    plt.rc("axes", titlesize=SIZE_LARGE)   # fontsize of the axes title
    plt.rc("axes", labelsize=SIZE_DEFAULT) # fontsize of the x and y labels
    plt.rc("xtick", labelsize=SIZE_DEFAULT)# fontsize of the tick labels
    plt.rc("ytick", labelsize=SIZE_DEFAULT)# fontsize of the tick labels

    fig, ax = plt.subplots(
        figsize=(6, 5)
    )

    # Generate a range of x values
    x_values = np.linspace(0, np.max(_data) + np.min(_data), 100)

    # Evaluate the function for each x value
    femp_y_values = [femp(x) for x in x_values]

    cdf_y_values = cdf_y_values(x_values)

    xticks = [i for i in range(0, int(np.max(_data) + np.min(_data)) + 1)]
    yticks = np.arange(0, 1.2 + 0.1, 0.1)

    # Hide the all but the bottom spines (axis lines)
    ax.spines["right"].set_visible(False)
    ax.spines["left"].set_visible(False)
    ax.spines["top"].set_visible(False)

    # Only show ticks on the left and bottom spines
    ax.yaxis.set_ticks_position("left")
    ax.xaxis.set_ticks_position("bottom")
    ax.spines["bottom"].set_bounds(min(xticks), max(xticks))

    # Plot femp(x)
    plt.plot(x_values, femp_y_values, label='femp(x)', color=RED)

    # Plot the cumulative distribution function
    plt.plot(x_values, cdf_y_values, label='CDF(x)', color='black')

    # plot y = 1 line
    plt.plot(x_values, np.full_like(x, 1), label='y = 1', 
             linestyle='--', color='black')

    # axis names
    plt.xlabel('x')
    plt.ylabel('F(x)')

    plt.xticks(xticks)
    plt.yticks(yticks)

    # Adjust the font size of the tick labels
    plt.tick_params(axis='both', which='major', labelsize=SIZE_TICKS)

    plt.grid(True)

    plt.show()
\end{lstlisting}

\begin{lstlisting}
# uniform
def uniform_y_values(x_values):
    return [sp.stats.uniform.cdf(
                x, 
                loc=np.min(_data), 
                scale=np.ptp(_data)
            ) for x in x_values]

buildFemp(uniform_y_values)
\end{lstlisting}

\begin{lstlisting}
# norm 
def norm_y_values(x_values):
    mean, std = sp.stats.norm.fit(data)
    return [sp.stats.norm.cdf(x, mean, std) for x in x_values]

buildFemp(norm_y_values)
\end{lstlisting}

\begin{lstlisting}
# beta 
def beta_y_values(x_values):
    a, b, loc, scale = sp.stats.beta.fit(data)
    return [sp.stats.beta.cdf(x, a, b, loc, scale) for x in x_values]

buildFemp(beta_y_values)
\end{lstlisting}

\end{document}
